\section*{Abstract}

Dermoscopy technique for classification of skin lesions may be useful in early
diagnosis of severe skin diseases and in the clinical management of skin
lesions. The challenge is to create an accurate and clinically usable system.
We create two versions of this design, based on requirements. The first one is
light in weight. It adopts a truncated MobileNetV2 backbone along with a compact
transformer, and with only 3.98 million parameters, it achieves 60.8\% BMA
(balanced multi-class accuracy). This compact design facilitates easy
installation within typical clinic environments, where space is a constraint.
The second version emphasizes more balanced accuracy than capacity of the
parameters. It is composed of two different pre-trained backbones
(DenseNet-121 and EfficientNet-B0) and a learned attention mechanism that
selectively fuses information from both, followed with a Kolmogorov--Arnold
network as the final classifier, and attains a BMA of 64.4\% on 16.45M
parameters. A dermatologist might consider the clinical context in the
classification process, but this information is seldom taken into account by
image-based classifiers. In this case, both models include patient metadata as
added diagnostic context directly into the classification pipeline, with the
data for age, sex and location of the body part where the lesion occurred. The
models provide an efficient and accurate balance for real-world dermatology
applications over same-budget baselines.

\vspace{1cm}
\textbf{Keywords:} Skin lesion classification; ISIC-2019; dermoscopy; hybrid
CNN--Transformer; Vision Transformer; clinical metadata fusion; cross-attention;
Kolmogorov--Arnold Network; class imbalance; balanced multi-class accuracy;
lightweight deep learning.
\pagebreak
